% =====================================================================
% FINAL PROJECT PROPOSAL -- Advanced NLP (CS7.501), Monsoon 2026
% Deadline 2.2: 20 August 2026.  Spec: 3-4 pages excluding references,
% covering project idea, motivation, background, related work, proposed
% methodology, dataset and evaluation plan, with an optional timeline.
% Submit as <TeamNameWithNoSpaces>-Proposal.pdf
% =====================================================================
\documentclass[11pt]{article}
\usepackage{acl}
\usepackage{times}
\usepackage{latexsym}
\usepackage[T1]{fontenc}
\usepackage[utf8]{inputenc}
\usepackage{microtype}
\usepackage{amsmath,amssymb}
\usepackage{booktabs}

\newcommand{\idf}{\mathrm{idf}}
\newcommand{\sfreq}{\mathrm{sf}}

\title{Predicting When Collection Statistics Help Top-$k$ Selection:\\
From Retrieval to Attention and Expert Routing}

\author{
  \textbf{Team ANLP Doomsday} \\[2pt]
  Aviral Gupta \quad Mohit Kumar Singh \quad Anirudh Sankar \quad Arihant Tripathy \\
  \texttt{\{aviral.gupta, mohit.singh, anirudh.sankar, arihant.tripathy\}@research.iiit.ac.in} \\
  IIIT Hyderabad, India
}

\begin{document}
\maketitle

\begin{abstract}
On sparse-autoencoder features, BM25 scores $0.582$ nDCG@10 where a plain inner
product over the same vectors scores $0.381$. On SPLADE the ordering reverses,
$0.455$ against $0.568$. What BM25 adds is collection frequency, the rate at
which a feature fires across a corpus, and the published account of the reversal
is that this helps only while a vocabulary stays Zipfian. That account has been
observed one vocabulary at a time and never made predictive, so for a new sparse
representation the only way to pick a scorer is to run both. We make the choice
predictable. Eight concept vocabularies built from one frozen encoder are scored
with a BM25-style function whose four components toggle independently, and we
regress the BM25-minus-inner-product margin on six distributional measurements
that cost one forward pass. Block-sparse attention and mixture-of-experts routing
pose the same choice, since both score a query against learned non-negative
features by inner product and neither consults selection frequency, so we apply
the fitted predictor to them without refitting. Whether it holds there is an open
question. Attention sinks fire for almost every query yet cannot be removed, and
the rarest experts are the most critical, so frequency may signal importance
differently inside a model than across a corpus. Evaluation runs on MS~MARCO, LIMIT and BEIR,
and on Hindi, Bengali, Telugu and Hindi-English code-mixed queries, where the
distributions the correction depends on are known to shift.
\end{abstract}

\section{Project Idea and Motivation}
\label{sec:idea}

A learned sparse retriever scores a query-document pair as
$\langle u_q, u_d \rangle$ over a concept vocabulary. A block-sparse attention
kernel estimates how critical KV page $i$ is for query $q_t$ as
$r_{t,i} = \langle q_t, c_i \rangle$ over page summaries, then attends to the
top-$k$ pages only \citep{tang2024quest}. A mixture-of-experts router gates on
$g_e = \langle h_t, w_e \rangle$ and sends the token to its top-$k$ experts
\citep{shazeer2017moe,fedus2022switch}. These belong to three different
literatures and perform one computation, an inner product against a learned
non-negative basis followed by a hard top-$k$.

What all three leave out is any account of how often a given feature wins.
Classical retrieval settled that question decades ago. The correction has three
parts, namely inverse document frequency, which discounts features that fire
everywhere, term-frequency saturation, so repeated evidence stops accumulating
linearly, and length normalisation, so that large documents are not rewarded
merely for being large \citep{robertson2009bm25}. Parts of this correction have
been reinvented inside the transformer under other names. Loss-free load
balancing adds a per-expert bias driven by recent load
\citep{wang2024lossfree}, which is an additive, unnormalised frequency discount
applied during training with no saturation or length term. The auxiliary losses
it replaces push toward uniform expert use, which is a cruder version of the
same idea \citep{zoph2022stmoe}. Streaming inference hand-codes an exception for
the highest-frequency keys in the cache \citep{xiao2024streamingllm}. Each of
these is a local fix. No shared account exists of when a collection-level
correction should help, and there is no way to tell before running the
experiment.

Our project turns that into something measurable. We hold representation
learning fixed, vary how a sparse basis is built and how it is scored, and ask
whether cheap measurements of a basis predict how much the correction helps.
The prediction is then tested in two places the measurements were never fitted
on, both inside the transformer.

\section{Background}
\label{sec:background}

BM25 scores a document $d$ against query $Q$ as
\begin{equation}
\sum_{i \in Q} \idf_i \cdot
\frac{f_i(d)\,(k_1+1)}{f_i(d) + k_1\!\left(1 - b + b\frac{|d|}{\overline{dl}}\right)},
\end{equation}
with $\idf_i = \ln\!\left(1 + \frac{N - df_i + 0.5}{df_i + 0.5}\right)$. None of
these ingredients requires words. Each needs only a representation with
non-negative coordinates, a meaningful per-dimension collection frequency $df_i$,
and a magnitude that admits a notion of length. Ordinary dense embeddings satisfy none of this. Coordinates are signed, every
dimension is active in every document, and there is no length analogue, which is
why dense retrieval is scored with a plain inner product instead.

Learned sparse features change the picture. Sparse-autoencoder activations,
$k$-means codes and quantised codebooks are non-negative, sparse and count-like,
The objects the transformer selects over are not themselves non-negative, since
KV page summaries are built from signed keys, but the quantity IDF actually needs
is a selection count, and that is non-negative by construction for both pages and
experts. A page or an expert has a well-defined selection rate over a calibration
set, and both admit a magnitude. All three settings are therefore eligible for
the same correction, and whether it helps is an empirical question.

\section{Related Work}
\label{sec:related}

\paragraph{Learned sparse retrieval and its vocabularies.} SPLADE and its
successors learn term weights over the backbone subword vocabulary and retain
the inverted index \citep{formal2021splade,lassance2024spladev3}, with COIL
taking a contextualised view of exact match \citep{gao2021coil} and a unified
framework since organising the design space \citep{nguyen2023unified}. A known
pathology of this family is that some terms become badly over-activated, which
degrades efficiency \citep{mackenzie2021wacky}. Recent work replaces the subword
vocabulary with a learned concept vocabulary. CL-SR indexes SAE-decomposed dense
embeddings as concept units \citep{park2025clsr}, SPLARE trains an SAE-based
sparse retriever directly \citep{formal2026splare}, and SAE-SPLADE swaps
SPLADE's MLM head for an SAE encoder \citep{zong2026tokens}. All three score
with an inner product.

\paragraph{Scoring learned sparse features with BM25.}
\citet{clavie2026latent} show that SAE features taken from a frozen dense
retriever have roughly Zipfian collection statistics and are directly
BM25-scorable. Their Appendix~C holds the sharpest evidence for our question.
On the same suite, BM25 beats the inner product on SAE latent terms ($0.582$
against $0.381$ nDCG@10) and loses to it on SPLADE-v3 ($0.455$ against
$0.568$), so the two scorers reverse order when the vocabulary changes beneath
them. \citet{han2026bm25v} run the comparison on image features and find
(their \S4.4) that the two differ by at most about two points of R@1,
attributing the parity to top-$k$ filtering already suppressing what IDF would
down-weight. The same section reports that raising the per-patch activation
budget from $k{=}16$ to $k{=}128$ makes document frequencies converge and drops
R@1 from $0.508$ to $0.038$. Both papers that ran the comparison propose the
same explanation, that BM25 helps while the vocabulary keeps its Zipfian
discriminative structure. That hypothesis is therefore already in print and we
are not supplying a missing explanation. It has only ever been observed on
individual vocabularies, never tested across a controlled family of them, never
quantified into anything usable before the fact, and never checked outside
retrieval.

\paragraph{Selection inside the transformer.} Quest scores KV pages by an inner
product against page min-max key summaries and loads only the top-$k$
\citep{tang2024quest}. H2O keeps a mix of recent and heavy-hitter tokens
\citep{zhang2023h2o}, and StreamingLLM preserves the earliest keys as attention
sinks \citep{xiao2024streamingllm}, which are entangled with the
massive-activation mechanism \citep{sun2024massive}. Expected Attention ranks a
cached pair by how strongly future queries are predicted to attend to it and
prunes the low scorers \citep{devoto2025expectedattention}. None of these uses
the rate at which a page is selected across queries. On the routing side,
\citet{su2026superexperts} show that a handful of rare experts are
disproportionately critical. To our knowledge no work treats retrieval scoring,
KV selection and routing as instances of one scoring problem, and none tries to
predict the effect of a correction before applying it.

\section{Proposed Methodology}
\label{sec:method}

\paragraph{A modular scorer.} For every setting we implement a single scoring
function in which soft IDF weighting, term-frequency saturation, length
normalisation and top-$k$ sparsification are separate switches. The uncorrected
inner product is the all-off configuration and BM25 is the all-on
configuration, so we can ask which mechanism is responsible for any gain rather
than only whether the bundle helps. The quantity we predict throughout is the
margin, meaning the change in task quality between two configurations at matched
budget.

\paragraph{Setting A: retrieval, used for calibration.} We build eight concept
bases from a single frozen encoder, spanning near-dense to strongly lexical, namely a
split-rectified identity map, random projection, PCA, $k$-means codes, product
quantisation \citep{jegou2011pq}, residual quantisation, a Top-$K$ sparse
autoencoder \citep{gao2025topksae}, and SPLADE
\citep{formal2021splade,lassance2024spladev3}. We include SPLADE on purpose,
since the reversal we want to explain is SPLADE against SAE and omitting it
would leave us unable to reproduce the phenomenon we study. Each basis is instantiated
at pooled and token-level granularity, giving sixteen representations, so
granularity becomes a variable rather than an unexamined confound.

\paragraph{Setting B: KV-page selection.} We start from Quest-style page
criticality \citep{tang2024quest} and add an online statistic $\sfreq_i$, the
fraction of prefill queries for which page $i$ enters the top-$k$. The corrected
score is
\begin{equation}
\tilde r_{t,i} = \phi(r_{t,i}) \cdot \idf(\sfreq_i) \cdot \ell(\|c_i\|),
\end{equation}
where $\phi$ is the saturation term and $\ell$ the length analogue. The method
is training-free and drops into an existing kernel. Whether attention sinks are
exempted from the correction is an ablation rather than an assumption. If IDF
suppresses them and accuracy falls, that outcome is the mechanism result.

\paragraph{Setting C: expert routing.} The same correction is applied to router
logits at inference time on OLMoE-1B-7B \citep{muennighoff2025olmoe}, which
activates 8 of 64 experts per layer and is open down to its training logs, with
expert selection frequency measured on held-out calibration data. Loss-free balancing
\citep{wang2024lossfree} is the additive, training-time special case we compare
against. We check directly whether the correction protects or destroys the
Super Experts identified by \citet{su2026superexperts}.

\paragraph{Diagnostics and the predictor.} Before any selection experiment, and
from a single forward pass over calibration data, we compute six measurements
per basis, namely active fraction per item, the Zipf exponent of the
selection-frequency curve, the Gini coefficient of activation mass, head mass,
the variance of IDF values, and sink share, meaning the fraction of total score
mass held by the top $1\%$ most-frequently-selected features. Prior work already
reports several of these descriptively, with \citet{han2026bm25v} fitting
power-law exponents in $[1.20, 2.32]$ and head fractions under $2\%$, but none
of it is used to predict anything. We regress the margin on these measurements
using the many cheap points Setting~A provides, validate by holding out whole
basis families, and then apply the fitted model without refitting to Settings~B
and~C. A predictor fitted on retrieval and tested inside a transformer is a
stronger claim than one cross-validated within a single grid, and it is the part
of the project a single-setting paper cannot produce.

\paragraph{Theory.} We aim to characterise the conditions under which an
IDF-weighted linear scorer over a learned basis provably reduces to a bilinear
inner product on the underlying dense embeddings. We call this the collapse
condition. Where it holds the two scorers must agree, which converts part of the
grid into something explained rather than tabulated.

\section{Datasets}
\label{sec:data}

Setting~A fits its bases and tunes hyperparameters on MS~MARCO passage
\citep{bajaj2016msmarco}, which supplies $8.8$M passages, and also uses LIMIT
\citep{weller2025limit},
which separates lexical from semantic signal by construction and so should make
any BM25 advantage easy to see. Configurations that look interesting transfer
zero-shot to a subset of BEIR \citep{thakur2021beir} and to TREC-DL 19/20, with
hyperparameters tuned only on MS~MARCO so that transfer tests generalisation
rather than additional fitting.

Because collection frequency is a property of a language and a corpus rather
than of a model, we extend Setting~A to the Hindi, Bengali and Telugu portions
of MIRACL \citep{zhang2023miracl} and to Hindi-English code-mixed queries, where
transliteration and mixing distort the document-frequency curve. Neither prior
SAE-based sparse retriever touches an Indic language, and the multilingual
evaluation in \citet{zong2026tokens} covers six MIRACL languages of which none
is Indian. Corpus-level measures of mixing already predict downstream behaviour
on such data \citep{kodali2022symcom}, and our diagnostics are measures of the
same kind, so this is a direct test of whether the fitted model captures a
property of the distribution or an artefact of English.

Setting~B uses RULER \citep{hsieh2024ruler}, needle-in-a-haystack and LongBench
on Llama-3.2-1B/3B-Instruct and Qwen2.5-7B-Instruct. Setting~C uses OLMoE-1B-7B
with WikiText for perplexity and MMLU and GSM8K for downstream accuracy.

\section{Evaluation Plan}
\label{sec:eval}

For Setting~A we report nDCG@10, MRR@10 and Recall@100/1000 alongside index
size, query-document FLOPs and latency, since a scoring function that improves
quality while destroying the efficiency benefit of sparsity would be a much less
useful finding. As a pipeline check we reproduce the Appendix~C reversal of
\citet{clavie2026latent} within a stated tolerance before trusting our own
numbers. Setting~B reports accuracy at matched token budget against Quest
\citep{tang2024quest} and H2O \citep{zhang2023h2o}, plus the KV budget needed
for a fixed accuracy; Setting~C reports perplexity, MMLU, GSM8K and routing
entropy.

The headline result is the diagnostics-to-margin regression, reported as
held-out $R^2$ under leave-one-basis-family-out validation, followed by the
cross-setting transfer error when the model fitted on Setting~A is applied to
Settings~B and~C. Component ablations switch each BM25 mechanism on and off to
identify which are necessary for which basis.

The grid produces several hundred pairwise comparisons, so we fix the analysis
in advance. The primary comparison is all-on against all-off per
representation; component ablations are secondary. Significance uses a paired
bootstrap over queries with Holm--Bonferroni correction within each family of
comparisons, and we report effect sizes alongside $p$-values rather than
$p$-values alone.

Confirming that BM25 helps on Zipfian vocabularies would only repeat what the
cited papers already predict, so we state in advance which outcomes would be
informative. The diagnostics might predict a large gain for KV-page selection
where the measured effect is negative, which would show contribution to the
output and discrimination between queries coming apart inside the transformer. A
basis family might have diagnostics saying BM25 should help where it measurably
does not, falsifying the Zipfian account as stated. A frequency-corrected router
might recover part of the auxiliary-loss benefit with no retraining.

\section{Timeline and Division of Labour}
\label{sec:timeline}

\begin{table}[t]
\centering\small
\begin{tabular}{@{}ll@{}}
\toprule
\textbf{Week} & \textbf{Milestone} \\
\midrule
1 & Frozen encoder; MS~MARCO index built \\
2--3 & Pooled grid; reproduce Appendix~C \\
4--5 & Token-level grid; MIRACL and code-mixed \\
6 & Setting~B built and run \\
7 & Setting~C routing \textbf{(Mid, 30 Sep)} \\
8 & Fit predictor; held-out-family validation \\
9 & Cross-setting transfer; collapse condition \\
10 & Ablations, significance \textbf{(Final, 31 Oct)} \\
\bottomrule
\end{tabular}
\caption{Ten-week plan, 20 August to 31 October 2026.}
\label{tab:schedule}
\end{table}

Gupta builds the retrieval infrastructure, the indexing pipeline and the eight
bases. Singh builds the modular scorer and the diagnostics, and fits the
predictor. Sankar takes Setting~B. Tripathy takes Setting~C and the collapse
condition, with Singh. Table~\ref{tab:schedule} gives the schedule against the
two remaining course deadlines.

The hardware available to us provides nodes of four GPUs with $11$\,GB of device
memory each, $128$\,GB of host RAM and several terabytes of local scratch. That
memory ceiling shapes the design. MS~MARCO is
encoded once
with a frozen DistilBERT-scale encoder, a job that shards cleanly across nodes
because passages are independent, and all eight pooled bases are then linear or
quantisation maps over the cached embeddings, so no re-encoding is needed.
Token-level granularity is capped at a $1$M-passage subsample, roughly $55$M
vectors stored fp16 as a memmap on node-local scratch, which fits the per-node
budget with room to spare. The MIRACL slices are far smaller than MS~MARCO.

Settings~B and~C are inference-only and every model is sized to fit inside a
single node. Llama-3.2-1B runs in fp16 on one card, while Llama-3.2-3B and
Qwen2.5-7B run $4$-bit across two, leaving headroom for a $32$K KV cache. We
chose OLMoE-1B-7B for Setting~C over a larger checkpoint such as Qwen3-30B-A3B
because its $6.9$B total parameters fit a node at $4$-bit where the $30$B model
would not, and the routing question we ask does not need the larger model, since
OLMoE already exposes 64 experts per layer. If time runs short we cut in this
order, token-level granularity for the quantisation bases first, then TREC-DL,
then the longest RULER context lengths.

\section{Risks and Limitations}
\label{sec:risks}

If the predictor is weak on held-out families, the project becomes a negative
result together with a resource. The diagnostics failing to predict the margin
is informative precisely because two published papers assert the mechanism, and
we would report the grid as a benchmark with the collapse condition as the
theoretical contribution. If Setting~C proves infeasible on the hardware we can
obtain we fall back to Settings~A and~B, which still support the transfer claim.
Our core citations are months old, so being scooped is a real risk; the
cross-setting transfer measurement is least exposed to it.

Four limitations are worth stating now. Setting~A uses a single frozen encoder,
so basis effects and encoder effects are not fully separable. Setting~C is
inference-time only, so we can make no claim about training dynamics. The
diagnostics are correlational, and the cross-setting transfer test is the only
part of the design with any causal flavour. Finally, the code-mixed evaluation
depends on query sets that are smaller and noisier than MS~MARCO, so those
results will carry wider confidence intervals than the English ones.

\bibliography{refs}

\end{document}
